\subsection{Jet coefficients and Boolean support}

The ideal-chain comparison has a second, coordinate-dependent companion. Fix a centered coordinate $t=s-x$ at a point of multiplicity $m\geq1$, and put
\[
 E_m=t\C[t]/(t^{m+1}),\qquad J_m=\C[t]/(t^m).
\]
The first is a pointed vector space, with ordered basis $t,\ldots,t^m$. Identify $\B^m$ with the subsets of $\{1,\ldots,m\}$ as in Theorem~\ref{thm:divisor}.

\begin{proposition}[The precise jet-support map]\label{prop:packet-jet-support}
The maps
\[
 \kappa_m\!\left(\sum_{r=1}^m c_rt^r\right)
       =\{r:c_r\ne0\},\qquad
 \iota_m(I)=\sum_{r\in I}t^r
\]
are maps of pointed sets and satisfy $\kappa_m\iota_m=\operatorname{id}$. The fibre of $\kappa_m$ over $I$ is $(\C^\times)^I$, with a singleton fibre over the empty set. The map $\kappa_m$ is not additive, and the additive monoids of $E_m$ and $\B^m$ are not isomorphic.

Division by the chosen coordinate gives a linear isomorphism $d_m:E_m\to J_m$. It takes $\C t^m$ onto the socle $\{a\in J_m:ta=0\}=\C t^{m-1}$. If a holomorphic germ $f$ has order $m$ at $x$, then
\[
 [f]=\frac{f^{(m)}(x)}{m!}t^m\ne0
       \quad\text{in }E_m.
\]
Its Boolean support is $\{m\}$, and $d_m[f]$ generates the socle of $J_m$.
\end{proposition}
\begin{proof}
The displayed basis gives unique coefficients, so support and the section are well defined. Prescribing the nonzero coefficients exactly at the indices in $I$ gives the stated fibre. However, $t$ and $-t$ both have support $\{1\}$, whereas their sum has empty support. Thus support fails to preserve addition. Every element of the target is additively idempotent, while $t+t\ne t$ in $E_m$, ruling out an additive-monoid isomorphism.

The map $d_m$ lowers every exponent by one and has inverse multiplication by $t$ followed by reduction modulo $t^{m+1}$. In $J_m$, multiplication by $t$ kills precisely the coefficient line of $t^{m-1}$. Finally the definition of vanishing order and Taylor expansion give
$f(s)=f^{(m)}(x)t^m/m!+O(t^{m+1})$, with nonzero leading coefficient. This proves all remaining assertions.
\end{proof}

For a closed discrete divisor, these maps assemble coordinatewise into maps of product sheaves. The support retraction is a retraction of pointed sheaves. It does not identify the vector-space jet sheaf with a Boolean semimodule sheaf. The order of vanishing and the socle line are intrinsic; the support of a general coefficient vector depends on the chosen coordinate. Divisor addition also has a direct Boolean interpretation: concatenation gives $\mathcal V_{D_1+D_2}\cong\mathcal V_{D_1}\times\mathcal V_{D_2}$. For disjoint supports this is simply the separation of the corresponding factors.

The Boolean coordinates also have their ordinary polynomial meaning. For every commutative $\B$-algebra $A$, evaluation gives
\[
 \operatorname{Hom}_{\B\text{-alg}}(\B[T_1,\ldots,T_m],A)
       \cong A^m.
\]
An image of each variable determines exactly one homomorphism, by the polynomial universal property. On $\B$-valued points, the substitutions $T_j\mapsto T$ and $T\mapsto T_1+\cdots+T_m$ induce the diagonal and join maps. Their composite sends $T$ to $mT=T$ in the idempotent polynomial semiring, so the retraction is already realized by these polynomial maps. This does not make the join map $\B^m\to\B$ a semiring homomorphism: for $m\geq2$, two disjoint nonempty coordinate supports have zero coordinatewise product but their joins both equal $1$.

An endpoint-label construction is equally explicit. In the two-dimensional space of principal parts
\[
 V=\C\frac{ds}{s}\oplus\C\frac{ds}{s-1},
\]
the residues at $0,1$ identify $V$ with $\C^2$. Subsets of $\{0,1\}$ correspond exactly to its coordinate subspaces. Thus the support map from $R_0^2$ to these four subspaces forgets precisely the coefficient fibre $R^I$ above the support $I$. If
\[
 q(s)=\frac1{s(s-1)},\qquad
 \ell(s)=\frac1s+\frac1{s-1},
\]
then $Aq(s)ds+B\ell(s)ds$ has residues $(-A+B,A+B)$; the inverse is $A=(b-a)/2$, $B=(a+b)/2$. Pullback by $s\mapsto1-s$ exchanges the endpoint basis forms. These are linear and support-coordinate identities, with no additional analytic information about a divisor.

\subsection{Descent and phase decorations}

The stabilizer qualification in the permutation construction can be stated explicitly.

\begin{proposition}[Invariant pushforward on a discrete orbit space]\label{prop:packet-stabilizer}
Let a finite group $G$ act on a discrete set $Z$, and suppose pointed sets or algebraic objects $M_z$ have compatible transport isomorphisms $M_z\to M_{gz}$. Let $q:Z\to Z/G$. For the product sheaf $\mathcal F$ and an orbit $O=Gz$,
\[
 (q_*\mathcal F)^G_O\cong M_z^{G_z},
 \qquad G_z=\{g:gz=z\}.
\]
For permutation-only transport on $\B^{m_z}$, with multiplicity constant on each orbit, this is the diagonal copy of $\B^{m_z}$. The same statement holds for the chain of ideals of $J_{m_z}$, and its two adjoints to the Boolean-cube inclusion descend coordinatewise.
\end{proposition}
\begin{proof}
The quotient of a discrete space is discrete, so the stalk over $O$ is the product of its fibres. A $G$-invariant tuple is determined by its $z$-coordinate, and invariance forces that coordinate to be fixed by $G_z$. Conversely, given $a\in M_z^{G_z}$, assign $ga$ to the coordinate $gz$. If $gz=hz$, then $h^{-1}g\in G_z$, so the assignment is independent of the representative. The two constructions are inverse and preserve every given operation. With permutation-only transport the stabilizers act trivially. The chain-cube maps depend only on $m_z$, so they commute with transport and restrict to invariants.
\end{proof}

This construction applies to any multiplicity-preserving symmetry of a divisor. For completed-zeta packets it uses the ordinary conjugation and functional-equation symmetries; it provides no additional assertion about zero locations.

There is also a precise pointed-set meaning for a phase decoration. Set $M=\{*\}\sqcup S^1$, pointed at $*$, and define
\[
 \pi(*)=0,\quad \pi(z)=1\ (z\in S^1),\qquad
 \iota(0)=*,\quad \iota(1)=1\in S^1.
\]
Then $\pi\iota=\operatorname{id}_{\B}$. On a closed discrete set, products give a split map of pointed sheaves from $\prod M$ to $\prod\B$, with fibre $S^1$ over an occupied Boolean coordinate. At the simple zeros of a fixed-coordinate holomorphic function, the values $f'(x)/|f'(x)|$ give a section. This assertion supplies a pointed decoration; it does not equip $M$ with an inverse-semiring module structure. In particular, analytic phase-product identities require their own analytic hypotheses and do not follow from the Boolean support map.

\subsection{The vector tail has no additional higher cohomology}

For completeness, the vector-space counterpart of Proposition~\ref{prop:tail} identifies the cohomology of the finite-support extension.

\begin{proposition}[Tail exact sequence]\label{prop:packet-tail-exact}
Let $Z\subset\C$ be closed and discrete, let $V_z$ be complex vector spaces, and let $\mathcal F(U)=\prod_{z\in U\cap Z}V_z$. Put
\[
 T=\left(\prod_{z\in Z}V_z\right)
       \Big/\left(\bigoplus_{z\in Z}V_z\right).
\]
For $j:\C\hookrightarrow\mathbf P^1(\C)$ and $i_\infty:\{\infty\}\hookrightarrow\mathbf P^1(\C)$ there is an exact sequence
\[
 0\longrightarrow j_!\mathcal F\longrightarrow j_*\mathcal F
   \longrightarrow i_{\infty*}T\longrightarrow0.
\]
Moreover $H^q(\mathbf P^1(\C),j_!\mathcal F)=0$ for every $q\geq1$.
\end{proposition}
\begin{proof}
On an open set containing infinity, the extension by zero consists of the finite-support families. Indeed its support must avoid some neighborhood of infinity, and a compact set meets $Z$ in finitely many points. On an open set avoiding infinity it is the unrestricted product sheaf. These rules give a sheaf: in a cover of a neighborhood of infinity, choose one member containing infinity; only finitely many support points lie outside that member, so gluing preserves finite support.

The map to the tail class on neighborhoods of infinity, and the zero map on other opens, is compatible with restriction. Its stalk kernel is $j_!\mathcal F$ by Proposition~\ref{prop:tail}; its stalk map at infinity is surjective, and all other target stalks vanish. This proves exactness. Both $j_*\mathcal F$ and $i_{\infty*}T$ are flasque, since their restriction maps are coordinate projections or identity and zero maps. Flasque sheaves are acyclic \citep[Lemma 20.12.3, Tag 09SY]{StacksFlasque}. Their long exact cohomology sequence gives the required vanishing in degrees at least two. In degree one it remains to observe that the global map $\prod V_z\to T$ is surjective by the definition of $T$. This proves the degree-one assertion as well.
\end{proof}
